\documentclass[11pt]{article}
\usepackage[margin=1in]{geometry}
\usepackage{amsmath,amssymb,amsthm}
\usepackage[hidelinks]{hyperref}
\setlength{\emergencystretch}{1em}
\newtheorem{proposition}{Proposition}
\newcommand{\R}{\mathbb{R}}
\newcommand{\C}{\mathbb{C}}
\title{Disproof of Conjecture 00000007668\\\large The geometric mean of finite q-Pochhammer zeros}
\author{C0ldSmi1e}
\date{October 5, 2026}
\begin{document}
\maketitle

\paragraph{The explicit source assertion.}
Both language versions define $g_N$ as the nonnegative $N$th root of the
product of the moduli of all $N$ zeros of $(z;q)_N$, and assert, for each
fixed real $0<q<1$, that
\[
 g_N=q^{-(N-1)/2}(1+q)^{-1}\bigl(1+O(q^{N/2})\bigr)
 \quad\text{as }N\longrightarrow\infty.
\]
Here $N$ is a positive integer, zeros are counted with algebraic multiplicity,
and all exponents and divisions are real. The implicit constant and starting
index in the $O$ term may depend on $q$.

\paragraph{Meaning of the polynomial.}
We use the standard finite q-Pochhammer notation \cite{DLMF}:
\[
 P_{q,N}(z)=(z;q)_N=\prod_{j=0}^{N-1}(1-zq^j),\qquad z\in\C.
\]
The source calls this a ``partial q-binomial'' without giving another formula.
The displayed standard symbol supplies the convention used here. An
alternative truncated series would require its own definition; it is not
silently substituted for this polynomial.

\begin{proposition}
For every real $0<q<1$ and every positive integer $N$, the geometric mean of
all complex zero moduli of $P_{q,N}$ is exactly $q^{-(N-1)/2}$.
Consequently the displayed relative-error assertion in the source is false.
\end{proposition}
\begin{proof}
Each factor is a nonzero degree-one polynomial whose root is $q^{-j}$.
The product has degree $N$, and its complete multiset of roots is therefore
$q^{-j}$ for $j=0,\ldots,N-1$, counting the multiplicities supplied by these
factors. All these numbers are positive real numbers. Hence
\[
 \prod_{i=1}^{N}|z_i|
 =\prod_{j=0}^{N-1}q^{-j}
 =q^{-N(N-1)/2},\qquad
 g_N=\left(q^{-N(N-1)/2}\right)^{1/N}=q^{-(N-1)/2}.
\]
Let $A_N=q^{-(N-1)/2}/(1+q)$. It is positive, and the source's relative error
is thus exactly
\[
 \frac{g_N}{A_N}-1=q.
\]
If the asserted $O$ bound held, there would be a real $C>0$ and an integer
$N_0\ge1$ such that $q\le Cq^{N/2}$ for every $N\ge N_0$.
But $q^{N/2}\to0$ for fixed $0<q<1$, so this would imply $q\le0$,
contradicting the hypothesis.
\end{proof}

\newpage
\raggedright
\paragraph{Relation to the whole conjecture.}
The displayed asymptotic formula is a necessary assertion in both language
versions. Its negation disproves the conjunction under the declared standard
notation. The later phrases do not specify the theta functions, Jensen mean,
correction coefficient or referent of ``this limit constant.'' This argument
does not assign definitions to those phrases or separately settle the stated
transcendence refinement. Establishing the displayed formula alone would
not prove those refinements; refuting it suffices for the stated conjunction.


\paragraph{Exact formal target.}
In \texttt{Conjecture7668.lean}, \texttt{qPochhammer} is a polynomial over
$\C$. Its evaluation is proved to equal the displayed finite product.
\texttt{rootModulusProduct} multiplies the norms in Mathlib's actual
\texttt{Polynomial.roots} multiset. The theorems
\texttt{qPochhammer\_roots} and \texttt{qPochhammer\_roots\_card}
prove both the full multiset identity and its cardinality $N$.
\texttt{geometricMean} uses \texttt{Real.rpow} with exponent $1/(N:\R)$;
nonnegativity and the $N$th-power identity are proved for $N>0$.
The totalized value at $N=0$ is irrelevant to the eventual assertion.

\texttt{RelativeAsymptotic q} is exactly
\[
 \exists C\in\R,\ C>0\ \land\ \exists N_0\in\mathbb N,\ N_0\ge1\ \land\
 \forall N\ge N_0,\quad
 \left|\frac{g_N}{A_N}-1\right|\le Cq^{N/2}.
\]
The lemma \texttt{relativeAsymptotic\_iff\_isBigO} proves equivalence
to Mathlib's \texttt{Asymptotics.IsBigO} along \texttt{Filter.atTop}.
There is no uniformity requirement in $q$, finite cutoff, or restriction to a
subsequence. The exponents use real subtraction and real division.

\texttt{NecessaryClause} universally quantifies the preceding statement
over all $0<q<1$. The theorem \texttt{conjecture7668\_disproof} proves its
negation; \texttt{not\_relativeAsymptotic} and
\texttt{not\_relative\_isBigO} give the stronger failure for every such $q$.
The final logical lemma \texttt{no\_completion} accepts the premise
$\mathrm{FullSource}\Rightarrow\mathrm{NecessaryClause}$ and derives
$\neg\mathrm{FullSource}$. This premise records the semantic implication
explained above; it is not a formal definition of the source's unspecified
later refinements.

\paragraph{Verification and reproduction.}
Lean 4.19.0 and Mathlib v4.19.0 are pinned by the included toolchain and
manifest. The unchanged proof was rebuilt in fresh independent projects,
then both proof and inspection sources were replayed with warnings treated
as errors. The audit covers all 29 authored declarations (8 definitions and
21 theorems), plus all 9 additional compiler-generated constants.
Their types, definitions, axiom dependencies and logical dependency
closures were inspected. Only \texttt{propext}, \texttt{Classical.choice}
and \texttt{Quot.sound} occur; there are no admitted proofs, custom axioms,
\texttt{native\_decide}, or unsafe/partial proof dependencies.
All nine dependency revisions and tracked sources were checked before and
after replay. No numerical experiment or external mathematical computation
is needed. The auxiliary programs inspect and reproduce the formal checks.

From the included \texttt{lean/} directory, after fetching the pinned
packages and Mathlib cache, run \texttt{lake build}, then
\texttt{lake env lean -DwarningAsError=true Conjecture7668.lean}
and the same command for \texttt{Check.lean}. The submission's
\texttt{VERIFICATION.md} gives the full audit commands, exact file hashes,
independent review and report checks. These are local validation records;
maintainer acceptance is a separate decision.

\begin{thebibliography}{9}
\raggedright
\bibitem{DLMF} NIST Digital Library of Mathematical Functions,
\href{https://dlmf.nist.gov/17.2.E1}{equation 17.2.1}, standard finite
q-Pochhammer symbol.
\end{thebibliography}
\end{document}
